\documentclass[10pt,a4paper]{amsart}
\usepackage[margin=19mm]{geometry}
\usepackage{amsmath,amssymb,mathtools}
\usepackage[scr=rsfs,cal=euler]{mathalfa}
\usepackage{xfrac}
\usepackage[unicode,hidelinks,bookmarks=false]{hyperref}
\newcommand{\ie}{i.e.\ }
\newcommand{\cf}{cf.\ }
\newcommand{\resp}{respectively\ }
\newcommand{\Subsection}{\subsection}
\providecommand{\nobreakdash}{\nobreak\hbox{-}}
\setlength{\emergencystretch}{2em}
\multlinegap0pt
\usepackage{bm}
\usepackage{bbm}
\usepackage{dsfont}
\usepackage{xspace}
\usepackage{cancel}
\usepackage{mathtools}
\usepackage[shortlabels]{enumitem}
\usepackage{amsthm}
\makeatletter
\@ifundefined{theorem}{\newtheorem{theorem}{Theorem}}{}
\@ifundefined{lemma}{\newtheorem{lemma}[theorem]{Lemma}}{}
\@ifundefined{proposition}{\newtheorem{proposition}[theorem]{Proposition}}{}
\makeatother
\DeclareMathOperator{\supp}{supp}
\newcommand{\E}{\mathbb{E}}
\newcommand{\Ind}{\mathbbm{1}}
\newcommand{\N}{\mathbb{N}}
\newcommand{\R}{\mathbb{R}}
\newcommand{\norm}[1]{\|{#1}\|}
\title[Optimal convex $M$-estimation via score matching]{Optimal convex $M$-estimation via score matching}
\author{Oliver Y. Feng, Yu-Chun Kao, Min Xu, Richard J. Samworth}
\date{Source: arXiv:2403.16688v2 (CC BY 4.0)}
\begin{document}
\maketitle

\noindent\emph{Source and licence.} Optimal convex $M$-estimation via score matching. Version arXiv:2403.16688v2 (CC BY 4.0), distributed under the Creative Commons Attribution 4.0 International licence, as recorded in the arXiv licence metadata for this exact version and shown on the abstract page https://arxiv.org/abs/2403.16688v2. Full rights notice: licenses/arxiv-2403.16688-CC-BY-4.0.txt.\par\smallskip

\noindent\textit{Editorial scope.} Counted unit: the Lemma whose source label is \texttt{lem:pk-limit} in the article (source0.tex lines 2040--2053, source label \texttt{lem:pk-limit}), delivered together with the article's own complete original proof of it (source0.tex lines 2055--2139, one \texttt{proof} environment of 85 lines). No mathematics is written, repaired or re-derived: statement and proof are the article's own text line for line. Required context retained uncounted: lem:psi0-star (lines 269--283); lem:p0-star (lines 430--433); prop:p0-star-fisher (lines 583--597); lem:J0-J0hat (lines 1553--1556); lem:p0-psi0-ineq (lines 1719--1738); lem:pk-induction (lines 1902--1917); eq:rk (lines 1936--1943). Every definition, display and label of those lines is retained unchanged. Omitted from this package and not redistributed: 1--268, 284--429, 434--582, 598--1552, 1557--1718, 1739--1901, 1918--1935, 1944--2039, 2054--2054, 2140--3904. Nothing inside a retained excerpt is omitted or altered: every retained body is delivered exactly as the article's lines are written, so no omission receipt of any kind accompanies this package. Citations: the retained excerpts carry no citation command at all, so no bibliography entry of the article is transcribed and no reference is redistributed. Reference adaptation: none was needed and none was made; every reference inside the delivered unit resolves inside it, so no unresolved reference or citation is produced. Printed slips kept literal: no printed word, symbol, hypothesis or number of the counted statement or of its proof was altered. Layout and class: the article's own class is \texttt{article} with packages including geometry, amsmath, amsthm, amssymb, bbm, hyperref, enumitem, graphicx, caption, subcaption, centernot, algorithm2e, natbib, xcolor; the wrapper is the lane's pinned \texttt{amsart} editorial wrapper at 10pt on A4, which restates the theorem-like environments the retained text needs and adds the article's own macro definitions that the retained text needs (\texttt{\textbackslash E}, \texttt{\textbackslash Ind}, \texttt{\textbackslash N}, \texttt{\textbackslash R}, \texttt{\textbackslash norm}, \texttt{\textbackslash supp}), with extra wrapper packages (bm, bbm, dsfont, xspace, cancel, mathtools, enumitem). The delivered numbering is the unit's own. Assets: no retained span contains a figure environment, a graphics-inclusion command, a PDF-page inclusion, a table or a TikZ picture; no image file is copied into this package, the package declares no asset dependency and no image file is redistributed with it. Licence: version arXiv:2403.16688v2 of this article is distributed under the Creative Commons Attribution 4.0 International licence, as recorded in the arXiv licence metadata for this exact version and shown on the abstract page https://arxiv.org/abs/2403.16688v2. A full rights notice is delivered at licenses/arxiv-2403.16688-CC-BY-4.0.txt. Characters: every retained excerpt is pure ASCII, so the delivered engine, XeLaTeX with its default Latin Modern text font, reports no missing character.
\begin{lemma}
\label{lem:psi0-star}
Let $P_0$ be a distribution with a uniformly continuous density $p_0$ on $\R$. Let $F_0 \colon [-\infty,\infty] \to [0,1]$ be the corresponding distribution function, and for $u \in [0,1]$, define
\[
F_0^{-1}(u) := \inf\{z \in [-\infty,\infty] : F_0(z) \geq u\} \quad\text{and}\quad J_0(u) := (p_0 \circ F_0^{-1})(u).
\] 
% $Q^0$ is right-continuous with $\lim_{v \nearrow u}Q^_0(v) = F_0^{-1}(u)$, so $J_0$ is right-continuous with left limits given by $J_0(u-) := \lim_{v \nearrow u}J_0(v) = (p_0 \circ F_0^{-1})(u)$ for $u \in (0,1]$.
Then both $J_0$ and its least concave majorant $\hat{J}_0$ on $[0,1]$ are continuous, with $p_0 = J_0 \circ F_0$ on $\R$, and
\[
\psi_0^* := \hat{J}_0^{(\mathrm{R})} \circ F_0
\]
is decreasing and right-continuous as a function from $\R$ to $[-\infty,\infty]$, provided that we set $\hat{J}_0^{(\mathrm{R})}(1) := \lim_{u \nearrow 1}\hat{J}_0^{(\mathrm{R})}(u)$.
% $= \hat{J}_0^{(\mathrm{L})}(1)$
Moreover, $\psi_0^*(z) \in \R$ if and only if $z \in \mathcal{S}_0$.
\end{lemma}

\begin{lemma}
\label{lem:p0-star}
In the setting of Lemma~\ref{lem:psi0-star}, there is a unique continuous log-concave density $p_0^*$ on $\R$ such that $\supp p_0^* = \mathcal{S}_0$ and $\log p_0^*$ has right derivative $\psi_0^*$ on $\mathcal{S}_0$. In particular, $\psi_0^* = (\log p_0^*)'$ Lebesgue almost everywhere on $\mathcal{S}_0$.  Furthermore, if $p_0$ is absolutely continuous, then $p_0^*$ minimises $I(p_0,p)$ over the class $\mathcal{P}_{\mathrm{LC}}$ of all univariate log-concave densities $p$, and if $p_0 \in \mathcal{P}_{\mathrm{LC}}$, then $p_0^* = p_0$.
\end{lemma} 

\begin{proposition}
\label{prop:p0-star-fisher}
For a uniformly continuous density $p_0 \colon \R \to \R$, the log-concave Fisher divergence projection $p_0^*$ and its corresponding distribution function $F_0^* \colon [-\infty,\infty] \to \R$ have the following properties.
\begin{enumerate}[label=(\alph*)]
\item Denote by $\mathcal{T}$ the set of $z \in \mathcal{S}_0$ such that $\psi_0^*$ is non-constant on every open interval containing~$z$. 
For $z \in \mathcal{T}$, we have
\begin{equation}
\label{eq:p0-star-hazard}
\frac{p_0^*(z)}{F_0^*(z)} \leq \frac{p_0(z)}{F_0(z)} \quad\text{and}\quad \frac{p_0^*(z)}{1 - F_0^*(z)} \leq \frac{p_0(z)}{1 - F_0(z)},
\end{equation}
whence $p_0^*(z) \leq p_0(z)$.
\item $\norm{p_0^*}_\infty \leq \norm{p_0}_\infty$ and $i(p_0^*) = -\int_\R p_0^*\,d\psi_0^* \leq -\int_\R p_0\,d\psi_0^* = i^*(p_0)$.
% Question/conjecture: Similarly to Corollary~\ref{cor:istar-cvx}, does this hold for any information quantity defined by minimising $V_{p_0}(\psi)$ over a convex cone of functions $\psi$?
\end{enumerate}
\end{proposition}

\begin{lemma}
\label{lem:J0-J0hat}
Let $\mathcal{T}$ be as in Proposition~\ref{prop:p0-star-fisher}\textit{(a)}. If $t \in \mathcal{T}$, then $p_0(t) = \hat{J}_0\bigl(F_0(t)\bigr) > 0$. 
\end{lemma}

\begin{lemma}
\label{lem:p0-psi0-ineq}
Let $s_0 := \inf\{z \in \R : \psi_0^*(z) \leq 0\}$ and $t_0 := \sup\{z \in \R : \psi_0^*(z) < 0\}$. For $s,t \in \mathcal{S}_0$ such that $s \leq t$, we have
\begin{equation}
\label{eq:psi0-int-integral}
\int_s^t \psi_0^* \,
\begin{cases}
\geq \log\hat{J}_0\bigl(F_0(t)\bigr) - \log\hat{J}_0\bigl(F_0(s)\bigr) \;\;&\text{if }t \leq t_0 \\
\leq \log\hat{J}_0\bigl(F_0(t)\bigr) - \log\hat{J}_0\bigl(F_0(s)\bigr)  \;\;&\text{if }s \geq s_0,
\end{cases}
\end{equation}
and hence
\begin{equation}
\label{eq:p0-psi0-ineq}
\begin{cases}
p_0(t) \leq p_0(s)\exp\bigl(\int_s^t \psi_0^*\bigr) \;&\text{if }t \leq t_0 \text{ and }s \in \mathcal{T} \\
p_0(t) \geq p_0(s)\exp\bigl(\int_s^t \psi_0^*\bigr) \;&\text{if }s \geq s_0 \text{ and }t \in \mathcal{T}.
\end{cases}
\end{equation}
\end{lemma}

\begin{lemma}
\label{lem:pk-induction}
For every $k \in \N_0$, the following statements hold.
\begin{enumerate}[label=(\alph*)]
\item $p_k \leq p_0$ on $\mathcal{T}$, and $\norm{p_k}_\infty = \norm{p_0}_\infty$;
\item $\displaystyle p_0(z) \int_s^t p_k \geq p_k(z) \int_s^t p_0$ if $z \in \{s,t\}$ for some $s,t \in \mathcal{T} \cup \{-\infty,\infty\}$ such that $s \leq t$;
% If either $t_k \leq t \wedge t_0$ or $s_k \geq s \vee s_0$
% NB: if $s_0 \in [s,t]$, then since $p_0(s_k) = p_k(s_0)$, we obtain the tighter inequality $\int_s^t p_k = \int_s^{s_0} p_k + \int_{s_0}^t p_k \geq \int_s^{s_0} p_{k-1} + \int_{s_0}^t p_{k-1} = \int_s^t p_{k-1}$.
\item $\dfrac{p_k(z)}{p_k(s_j)} = 
\begin{cases}
p_0(z)/p_0(s_j) &\text{if }z \in [s_j,t_j]\text{ for some }j \in \N\text{ with }j > k \\
p_j(z)/p_j(s_j) &\text{if }z \in [s_j,t_j]\text{ for some }j \in \N\text{ with }j < k;
\end{cases}$
\item $\displaystyle \psi_0^*(t-) \int_s^t p_k \leq p_k(t) - p_k(s) \leq \psi_0^*(s)\int_s^t p_k$ for all $s,t \in \mathcal{T} \cup \{-\infty,\infty\}$ such that $s \leq t$.
\end{enumerate}
\end{lemma}

\begin{equation}
\label{eq:rk}
r_k := \frac{p_{k-1}(t_k)e^{a_k(s_k - t_k)}}{p_{k-1}(s_k)} = \frac{p_0(t_k)e^{a_k(s_k - t_k)}}{p_0(s_k)}
\begin{cases}
\leq 1 \;\;&\text{if }s_k \leq s_0 \\
\geq 1 \;\;&\text{if }s_k \geq s_0.
\end{cases}
\end{equation}

\begin{lemma}
\label{lem:pk-limit}
For $z \in \R$, we have
\begin{equation}
\label{eq:pk-limit}
\lim_{k \to \infty}p_k(z) = \frac{p_0(s_0)}{p_0^*(s_0)} \cdot p_0^*(z).
\end{equation}
Moreover, let $t_{\min} := \inf\mathcal{T}$ and $t_{\max} := \sup\mathcal{T}$. If $s,t \in \mathcal{T} \cup \{-\infty,\infty\}$ are such that $-\infty < s \vee t_{\min} \leq t \wedge t_{\max} < \infty$, then
\begin{equation}
\label{eq:pk-int-limit}
\lim_{k \to \infty}\int_s^t p_k = \frac{p_0^*(s_0)}{p_0(s_0)} \int_s^t p_0^*.
\end{equation}
% Is this true in general if $|s| \vee |t| = \infty$?
\end{lemma}

\begin{proof}
For $k \in \N$, let $\phi_k := \log p_k$ and define $r_k$ as in~\eqref{eq:rk}. Then as noted above, $\phi_k > -\infty$ on $\mathcal{T}$, and Lemma~\ref{lem:pk-induction}\textit{(c)} ensures that
\begin{equation}
\label{eq:log-rk}
\log r_k = \psi_0^*(s_k)(t_k - s_k) - \bigl(\phi_{k-1}(t_k) - \phi_{k-1}(s_k)\bigr) = \int_{s_k}^{t_k}\psi_0^*(z)\,dz - \bigl(\phi_0(t_k) - \phi_0(s_k)\bigr).
\end{equation}
Fix $t \in \mathcal{T}$ such that $t \geq s_0$, and for convenience, define $[K] = \emptyset$ when $K = 0$ and $[K] = \N$ when $K = \infty$. Let $\mathcal{K}$ be the set of $\ell \in [K]$ such that $(s_\ell,t_\ell) \subseteq (s_0,t)$, so that $(s_\ell,t_\ell) \cap (s_0,t) = \emptyset$ for $\ell \in [K] \setminus \mathcal{K}$ and hence $\mathcal{T}^c \cap (s_0,t) = \bigcup_{\ell \in \mathcal{K}}(s_\ell,t_\ell)$. Similarly to the proof of Lemma~\ref{lem:p0-psi0-ineq}, let $U \sim U(0,1)$, so that $Z := F_0^{-1}(U)\sim P_0$, $F_0(Z) = U$ and $\psi_0^*(Z) = \hat{J}_0^{(\mathrm{R})}(U)$. We have $\psi_0^* \leq 0$ and $p_0 > 0$ on $\mathcal{T} \cap (s_0,t)$, and moreover $\hat{J}_0$ is Lipschitz on $(F_0(s_0),F_0(t)]$ with $\hat{J}_0^{(\mathrm{R})} \leq 0$, so
\begin{align}
\int_{s_0}^t \psi_0^* \Ind_{\mathcal{T}} &= \E\biggl(\frac{\psi_0^*(Z)}{p_0(Z)}\Ind_{\{Z \in \mathcal{T} \cap (s_0,t)\}}\biggr) = \E\biggl(\frac{\hat{J}_0^{(\mathrm{R})}(U)}{\hat{J}_0(U)}\Ind_{\{F_0^{-1}(U) \in \mathcal{T} \cap (s_0,t)\}}\biggr) \notag \\
&= \E\biggl(\frac{\hat{J}_0^{(\mathrm{R})}(U)}{\hat{J}_0(U)}\Ind_{\{F_0^{-1}(U) \in (s_0,t]\}}\biggr) - \sum_{\ell \in \mathcal{K}} \E\biggl(\frac{\hat{J}_0^{(\mathrm{R})}(U)}{\hat{J}_0(U)}\Ind_{\{F_0^{-1}(U) \in (s_\ell,t_\ell]\}}\biggr) \notag \\
% &= \E\biggl(\frac{\hat{J}_0^{(\mathrm{R})}(U)}{\hat{J}_0(U)}\Ind_{\{U \in (F_0(s_0),F_0(t)]\}}\biggr) - \sum_{k \in \mathcal{K}} \E\biggl(\frac{\hat{J}_0^{(\mathrm{R})}(U)}{\hat{J}_0(U)}\Ind_{\{U \in (F_0(s_0),F_0(t)]\}}\biggr) \notag \\
&= \int_{F_0(s_0)}^{F_0(t)} \frac{\hat{J}_0^{(\mathrm{R})}}{\hat{J}_0} - \sum_{\ell \in \mathcal{K}} \int_{F_0(s_\ell)}^{F_0(t_\ell)} \frac{\hat{J}_0^{(\mathrm{R})}}{\hat{J}_0} \notag \\
&= \log\hat{J}_0\bigl(F_0(t)\bigr) - \log\hat{J}_0\bigl(F_0(s_0)\bigr) - \sum_{\ell \in \mathcal{K}} \bigl\{\log\hat{J}_0\bigl(F_0(t_\ell)\bigr) - \log\hat{J}_0\bigl(F_0(s_\ell)\bigr)\bigr\} \notag \\
\label{eq:psi0-T}
&= \phi_0(t) - \phi_0(s_0) - \sum_{\ell \in \mathcal{K}}\bigl(\phi_0(t_\ell) - \phi_0(s_\ell)\bigr),
\end{align}
where the second and final equalities above follow from Lemma~\ref{lem:J0-J0hat} and the fact that $s_0,t,s_\ell,t_\ell \in \mathcal{T} \cup \{-\infty,\infty\}$ for $\ell \in [K]$. In addition, for all such $\ell$, we have $\phi_\ell(t) = \phi_{\ell-1}(t) + (\log r_\ell)\Ind_{\{\ell \in \mathcal{K}\}}$ by the definition of $p_\ell$, and $\phi_0^* := \log p_0^*$ satisfies $\phi_0^*(t) - \phi_0^*(s_0) = \int_{s_0}^t \psi_0^*$. Thus, by induction together with~\eqref{eq:log-rk} and~\eqref{eq:psi0-T},
\begin{align*}
\phi_k(t) - \phi_0(s_0) &= \phi_0(t) - \phi_0(s_0) + \sum_{\ell \in \mathcal{K} \cap [k]}\log r_\ell = \phi_0(t) - \phi_0(s_0) - \sum_{\ell \in \mathcal{K} \cap [k]} \Bigl(\phi_0(t_\ell) - \phi_0(s_\ell) - \int_{s_\ell}^{t_\ell}\psi_0^*\Bigr) \\
&= \biggl(\int_{s_0}^t \psi_0^* \Ind_{\mathcal{T}} + \sum_{\ell \in \mathcal{K}} \int_{s_\ell}^{t_\ell} \psi_0^*\biggr) + \sum_{\ell \in \mathcal{K} : \ell > k} \Bigl(\phi_0(t_\ell) - \phi_0(s_\ell) - \int_{s_\ell}^{t_\ell}\psi_0^*\Bigr) \\
&= \int_{s_0}^t \psi_0^* + \sum_{\ell \in \mathcal{K} : \ell > k} \Bigl(\phi_0(t_\ell) - \phi_0(s_\ell) - \int_{s_\ell}^{t_\ell}\psi_0^*\Bigr) \to \phi_0^*(t) - \phi_0^*(s_0)
% \begin{cases}
% = \phi_0^*(t) - \phi_0^*(s_0) &\text{if }k = K < \infty \\[3pt]
% \to \phi_0^*(t) - \phi_0^*(s_0) &\text{as }k \to \infty \text{ if }K = \infty,
% \end{cases}
\end{align*}
as $k \to \infty$, which yields~\eqref{eq:pk-limit} for $z = t$. On the other hand, if $t \in \mathcal{T}$ and $t < s_0$, then we can instead let $\mathcal{K}$ be the set of $\ell \in [K]$ such that $(s_\ell,t_\ell) \subseteq (t,s_0)$, and deduce by similar reasoning that
\begin{align*}
\phi_k(t) - \phi_0(s_0) &= \phi_0(t) - \phi_0(s_0) + \sum_{\ell \in \mathcal{K} \cap [k]} \Bigl(\phi_0(t_\ell) - \phi_0(s_\ell) - \int_{s_\ell}^{t_\ell}\psi_0^*\Bigr) \\
&= -\int_t^{s_0}\psi_0^* - \sum_{\ell \in \mathcal{K} : \ell > k} \Bigl(\phi_0(t_\ell) - \phi_0(s_\ell) - \int_{s_\ell}^{t_\ell}\psi_0^*\Bigr) \to \int_{s_0}^t \psi_0^* = \phi_0^*(t) - \phi_0^*(s_0)
% \begin{cases}
% = \phi_0^*(t) - \phi_0^*(s_0) &\text{if }k = K < \infty \\[3pt]
% \to \phi_0^*(t) - \phi_0^*(s_0) &\text{as }k \to \infty \text{ if }K = \infty,
% \end{cases}
\end{align*}
as $k \to \infty$. Having proved that~\eqref{eq:pk-limit} holds for all $z \in \mathcal{T}$, we now consider $z \in \mathcal{T}^c$, for which there exists a unique $\ell \in [K]$ such that $z \in (s_\ell,t_\ell)$. If $s_\ell > s_0$, then $p_\ell(z) = p_\ell(s_\ell)e^{\psi_0^*(s_\ell)(z - s_\ell)}$, while if $s_\ell \leq s_0$, then $p_\ell(z) = p_\ell(t_\ell)e^{\psi_0^*(s_\ell)(z - t_\ell)}$.
% \[
% p_\ell(z) = 
% \begin{cases}
% p_\ell(s_\ell)e^{\psi_0^*(s_\ell)(z - s_\ell)} &\text{if }s_\ell > s_0 \\
% p_\ell(t_\ell)e^{\psi_0^*(s_\ell)(z - t_\ell)} &\text{if }s_\ell \leq s_0.
% \end{cases}
% \]
For $k > \ell$, it follows from Lemma~\ref{lem:pk-induction}\textit{(c)} that
\[
p_k(z) = 
\begin{cases}
p_k(t_\ell)e^{\psi_0^*(s_\ell)(z - t_\ell)} &\text{if }s_\ell \leq s_0 \\
p_k(s_\ell)e^{\psi_0^*(s_\ell)(z - s_\ell)} &\text{if }s_\ell > s_0.
\end{cases}
\]
Moreover, $s_\ell,t_\ell \in \mathcal{T} \cup \{-\infty,\infty\}$, so if $s_\ell \leq s_0$, then
\[
p_k(z) = p_k(t_\ell)e^{\psi_0^*(s_\ell)(z - t_\ell)} \to \frac{p_0(s_0)}{p_0^*(s_0)} \cdot p_0^*(t_\ell)e^{\psi_0^*(s_\ell)(z - t_\ell)} = \frac{p_0(s_0)}{p_0^*(s_0)} \cdot p_0^*(z)
\]
as $K \to \infty$, and~\eqref{eq:pk-limit} holds similarly when $s_\ell > s_0$.
% \[
% \frac{p_0(s_0)}{p_0^*(s_0)} \cdot p_0^*(z) = \frac{p_0(s_0)}{p_0^*(s_0)} \cdot p_0^*(s_\ell)e^{\psi_0^*(s_\ell)(z - t_\ell)} =
% \begin{cases}
% p_K(s_\ell)e^{\psi_0^*(s_\ell)(z - t_\ell)} = p_K(z) &\text{if }K < \infty \\
% \lim_{k \to K}p_k(s_\ell)e^{\psi_0^*(s_\ell)(z - t_\ell)} = \lim_{k \to K}p_k(z) &\text{if }K = \infty,
% \end{cases}
% \]
% as required.

Finally, let $s,t \in \mathcal{T} \cup \{-\infty,\infty\}$ be such that $-\infty < s \vee t_{\min} \leq t \wedge t_{\max} < \infty$. If $t > t_{\max}$, then there exists a unique $j \in [K]$ such that $t_{\max} = s_j < t_j = t = \infty$, and by Lemma~\ref{lem:p0-star} and its proof, $\psi_0^*(t_{\max}) = \lim_{z \to \infty}\psi_0^*(z) \in (-\infty,0)$. Moreover, by Lemma~\ref{lem:pk-induction}\textit{(c)},
\[
p_k(z) = p_k(t_{\max})e^{\psi_0^*(t_{\max})(z - t_{\max})}
\]
for all $z > t_{\max}$ and $k \geq j$. Similarly, if $s < t_{\min}$, then there exists a unique $\ell \in [K]$ such that $-\infty = s = s_\ell < t_\ell = t_{\min}$, and $\psi_0^*(t_{\min}-) = \lim_{z \to -\infty}\psi_0^*(z) \in (0,\infty)$. In addition, by Lemma~\ref{lem:pk-induction}\textit{(c)},
\[
p_k(z) = p_k(t_{\min})e^{\psi_0^*(t_{\min}-)(z - t_{\min})}
\]
for all $z < t_{\min}$ and $k \geq \ell$. Thus, for $k \geq j \vee \ell$, it follows from Lemma~\ref{lem:pk-induction}\textit{(a)} that in all cases,
% $p_k(t) \leq p_0(t)$ for $t \in \{t_{\min},t_{\max}\} \subseteq \mathcal{T}$, so
\[
p_k(z) \leq
\begin{cases}
p_0(t_{\min})e^{\psi_0^*(t_{\min}-)(z - t_{\min})} \;\;&\quad\text{for }z \in (s,s \vee t_{\min}) \\
\norm{p_0}_\infty \;\;&\quad\text{for }z \in [s \vee t_{\min},t \wedge t_{\max}] \\
p_0(t_{\max})e^{\psi_0^*(t_{\max})(z - t_{\max})} \;\;&\quad\text{for }z \in (t \wedge t_{\max},t),
\end{cases}
\]
so the pointwise supremum $\sup_{k \geq j \vee \ell} p_k$ is integrable on $(s,t)$, and hence~\eqref{eq:pk-int-limit} follows from~\eqref{eq:pk-limit} and the dominated convergence theorem.
\end{proof}

\end{document}
